%% kapa: exam=1 round=29 year=2018 period=1 form=A subject=real_estate subject_ko=부동산학원론 seq=9 number=89
%% source: Q-Net 감정평가사 기출문제 (cst00309) artlSeq=5202808 file=extracted/1008/5202808/2018년 제29회 감정평가사 1차시험 기출문제/2018년 제29회 감정평가사 1차시험 기출문제/1교시 A형.pdf
%% conversion: kordoc 4.24.51 -> tools/gen1.py (auto); review=proofread-sonnet; flags=none
\begin{question}{89}
부동산정책의 시행으로 A지역 아파트시장의 공급함수는 일정하고 수요함수는 다음과 같이 변화되었다. 이 경우 y축, 수요곡선, 공급곡선으로 둘러싸인도형의 면적과 균형거래량의 변화는? (단, 거래량과 도형 면적의 단위는 무시하며, x축은 수량, y축은 가격을 나타냄)
\begin{framedtext}
○ 수요함수: $Q_{d1}=50-p$ (이전) → $Q_{d2}=80-p$ (이후) \par
○ 공급함수: $Q_{s}=-40+2p$ \par
○ $p$는 가격, $Q_{d}$는 수요량, $Q_{s}$는 공급량 \par
\end{framedtext}
\begin{choices}
  \item 면적: 700 증가, 균형거래량: 10 증가
  \item 면적: 900 증가, 균형거래량: 10 증가
  \item 면적: 700 증가, 균형거래량: 20 증가
  \item 면적: 900 증가, 균형거래량: 20 증가
  \item 면적: 700 증가, 균형거래량: 30 증가
\end{choices}
\end{question}
\answer{4}
